%% ============================================================
%%  サンプル付録
%%  本編から参照される前提知識をここに置く。
%%  \appendix 以降に \input するので、番号は A.1.1 の形になる。
%% ============================================================

\chapter{集合と写像}
\label{ch:preliminaries}

本編で用いる記法をまとめる。内容は標準的なもので、証明は省略する。

\section{集合}
\label{sec:sets}

\begin{definition}{冪集合}{power-set}
集合 $A$ の部分集合すべてからなる集合を $A$ の\textbf{冪集合}といい、$2^A$ と書く。
\end{definition}

\begin{definition}{上限}{supremum}
$S \subset \R$ を空でない上に有界な集合とする。$S$ の上界のうち最小のものを
$S$ の\textbf{上限}といい、$\sup S$ と書く。
\end{definition}

\begin{remark}{上限の存在}{supremum-exists}
空でない上に有界な $S \subset \R$ に対して $\sup S$ が存在することは、
実数の連続性公理そのものである。本資料ではこれを認める。
\end{remark}

\section{写像}
\label{sec:maps}

\begin{definition}{像と逆像}{image}
写像 $f \colon A \to B$ と $S \subset A$、$T \subset B$ に対し、
\[
  f(S) = \{ f(x) : x \in S \},
  \qquad
  f^{-1}(T) = \{ x \in A : f(x) \in T \}
\]
をそれぞれ $S$ の\textbf{像}、$T$ の\textbf{逆像}という。
\end{definition}

\begin{proposition}{逆像は演算を保つ}{preimage-operations}
写像 $f \colon A \to B$ と $T_1, T_2 \subset B$ に対して次が成り立つ。
\begin{itemize}
  \item $f^{-1}(T_1 \cup T_2) = f^{-1}(T_1) \cup f^{-1}(T_2)$
  \item $f^{-1}(T_1 \cap T_2) = f^{-1}(T_1) \cap f^{-1}(T_2)$
  \item $f^{-1}(B \setminus T_1) = A \setminus f^{-1}(T_1)$
\end{itemize}
像 $f(\cdot)$ については、和集合は保つが共通部分は一般に保たない。
\end{proposition}

\begin{proof}
いずれも $x \in f^{-1}(T)$ と $f(x) \in T$ が定義により同値であることから直ちに従う。
\end{proof}

\begin{example}{像は共通部分を保たない}{image-intersection}
$A = \{1, 2\}$、$B = \{0\}$、$f(1) = f(2) = 0$ とする。
$S_1 = \{1\}$、$S_2 = \{2\}$ とおくと $S_1 \cap S_2 = \emptyset$ だが、
$f(S_1) \cap f(S_2) = \{0\} \neq \emptyset$ である。
\end{example}
